\mypara{Visual concepts associated with human preference}
We apply frozen FG-CLIP2 dense image features and the full supplied concept bank to human-ranked alternatives within each dataset. For every concept, we average preferred-minus-rejected maximum patch similarities within annotation events, then within normalized prompts, then equally across prompts. ImageReward uses all strict ranked pairs; the linked four-candidate Pick-a-Pic release uses the selected image against each non-selected alternative. We freeze concepts and directions on discovery data and evaluate held-out prompts, preserving prompt and image dependencies. We report paired concordance and component-bootstrap uncertainty, with best-versus-rest and generator/patch-count sensitivities. These are associations with recorded choices, not calibrated prevalence estimates or causal explanations. This engineering pilot retained 15 held-out prompts in 15 dependence components for imagereward. Table entries report observed effects and per-concept intervals; visual interpretations await human review.
The first discovery-ranked terms were emerald-green, red-flag, red-velvet. Of 20 frozen candidates, 7 retained their discovery direction in held-out effects; this descriptive count is not a multiplicity-adjusted significance claim. The equal signed mean of the selected concepts achieved held-out paired concordance 0.453 (95\% component-bootstrap interval [0.352, 0.567]). These pilot findings remain provisional: visual concept presence is unverified, the supplied text-bank encoding provenance is unresolved, and the sample is small.
The linked four-candidate Pick-a-Pic metadata contained 25,355 ranking events, of which 13,085 met the strict metadata eligibility rules. However, all 119 image URLs attempted in the bounded smoke cohort returned HTTP 403. Consequently, Pick-a-Pic empirical effects and cross-dataset replication are unavailable; no alternative dataset release or model-generated preference labels were substituted.
